\documentclass[11pt,a4paper]{article}
\usepackage[utf8]{inputenc}
\usepackage[T1]{fontenc}
% babel-spanish opcional: descomente si tiene instalado el paquete
% \usepackage[spanish,es-noquoting]{babel}
\usepackage[margin=2.4cm]{geometry}
\usepackage{booktabs}
\usepackage{longtable}
\usepackage{fancyhdr}
\usepackage{xcolor}
\usepackage[hidelinks]{hyperref}
\usepackage{parskip}

\definecolor{uteqverde}{RGB}{27,94,60}

\pagestyle{fancy}
\fancyhf{}
\fancyhead[L]{\footnotesize UTEQ · Ingeniería de Requerimientos (ISR-401 / 20303)}
\fancyhead[R]{\footnotesize Desviaciones del pre-registro · CarniCore}
\fancyfoot[C]{\footnotesize Página \thepage}
\renewcommand{\headrulewidth}{0.4pt}

\begin{document}

\begin{center}
{\Large\bfseries\color{uteqverde}
Desviaciones respecto del protocolo pre-registrado}\\[4pt]
{\large Componente empírico · Sistema CarniCore · Enfoque 2}\\[8pt]
{\small Registro OSF: \url{https://osf.io/yp7t3} · Registrado el 2 de agosto de 2026}\\
{\small Protocolo de referencia: \texttt{06\_Experimento/protocolo.pdf}, versión 1.0}\\
{\small Ubicación en el repositorio: \texttt{06\_Experimento/osf\_deviations.pdf}}\\[4pt]
{\small Versión 1.0 · Elaborado para la entrega del examen final (semana 19)}
\end{center}

\vspace{6pt}
\hrule
\vspace{10pt}

\section*{1. Propósito y alcance}

La Sección 7.3 de la guía de la Entrega 4 (2B) exige documentar cada desviación
del análisis efectivamente ejecutado respecto del plan pre-registrado en OSF,
con la razón, el momento en que se detectó y la mitigación aplicada. Este
documento cumple esa exigencia y se replica en la sección \textit{Deviations
from pre-registration} del propio registro OSF.

Actualizar el registro con las desviaciones no invalida el pre-registro. Al
contrario: es la práctica que el uso maduro del pre-registro
exige\footnote{Nosek, B.~A., Ebersole, C.~R., DeHaven, A.~C. y Mellor, D.~T.
(2018). The preregistration revolution. \textit{PNAS}, 115(11), 2600--2606.}.
Lo que sí invalidaría el trabajo sería ejecutar un análisis distinto del
registrado y no declararlo.

\section*{2. Resumen ejecutivo}

El plan de análisis pre-registrado se ejecutó \textbf{íntegramente}. Se
registran dos desviaciones menores (DEV-01 y DEV-02) y una corrección
documental de gravedad mayor (COR-01) relativa a cómo se reportaron los
resultados en una versión preliminar del manuscrito.

\begin{longtable}{p{1.6cm}p{5.2cm}p{2.2cm}p{5.6cm}}
\toprule
\textbf{ID} & \textbf{Naturaleza} & \textbf{Detectada} & \textbf{Efecto sobre las
conclusiones} \\
\midrule
\endhead
DEV-01 & Ampliación del corpus de 25 a 27 RF & 21/08/2026 & Ninguno; extensión
aditiva verificable \\
DEV-02 & Ubicación de las notas de desacuerdo & 01/09/2026 & Ninguno; material
suplementario \\
COR-01 & Reporte erróneo en el manuscrito preliminar & 03/09/2026 & Corregido;
ver Sección 5 \\
DEV-03 & Corpus \texttt{rf27.json} desactualizado respecto del ERS v2.0 & 03/09/2026
(detectado) · 14/09/2026 (corregido) & Ninguno sobre la conclusión; ver Sección 5 bis \\
\bottomrule
\end{longtable}

\section*{3. DEV-01 --- Ampliación del corpus de 25 a 27 requisitos}

\textbf{Qué se registró.} El protocolo versión 1.0 (Sección 7, Población)
declara como población los 25 requisitos funcionales del ERS/SRS vigente al
2 de agosto de 2026.

\textbf{Qué se ejecutó.} El análisis se aplicó sobre 27 requisitos.

\textbf{Razón.} El comité de control de cambios del proyecto aprobó dos
requisitos adicionales con posterioridad al registro: RF-26 (vista pública de
trazabilidad por código QR con campos delimitados) y RF-27 (filtrado de
reportes gerenciales por proveedor). Excluirlos habría analizado un corpus que
ya no correspondía a la especificación vigente.

\textbf{Momento de detección.} 21 de agosto de 2026, al cerrar la versión 2.0
del ERS/SRS.

\textbf{Mitigación.} La extensión es puramente aditiva: los 25 requisitos
originales se conservan sin alteración alguna. Esto es verificable de forma
mecánica comparando los archivos \texttt{rf25.json} y \texttt{rf27.json} del
paquete de replicación; los 25 primeros elementos son idénticos campo a campo.
Las reglas del detector no se modificaron en ningún momento después del
registro.

\textbf{Efecto sobre las conclusiones.} Ninguno. El detector no marcó ningún
requisito ni entre los 25 originales ni entre los 2 añadidos.

\section*{4. DEV-02 --- Ubicación de las notas interpretativas de desacuerdo}

\textbf{Qué se registró.} La Sección 9, paso 6 del protocolo prevé redactar una
nota interpretativa breve por cada requisito en que el detector y el consenso
experto difieran.

\textbf{Qué se ejecutó.} Las cuatro notas (RF-08, RF-17, RF-21 y RF-22) se
redactaron, pero se publican como material suplementario en el paquete Zenodo y
no en el cuerpo del artículo.

\textbf{Razón.} Restricción de extensión del formato de la conferencia
objetivo.

\textbf{Momento de detección.} 1 de septiembre de 2026, al maquetar el
manuscrito.

\textbf{Mitigación.} Las notas se incluyen íntegras en el paquete de
replicación y se referencian explícitamente desde la sección de resultados del
manuscrito.

\textbf{Efecto sobre las conclusiones.} Ninguno.

\section*{5. COR-01 --- Corrección de un reporte erróneo en el manuscrito
preliminar}

Este apartado no documenta una desviación del plan de análisis, sino un error
de reporte que se detectó y se corrigió. Se declara con el mismo detalle por
transparencia.

\textbf{Qué ocurrió.} La versión del manuscrito circulada el 1 de septiembre de
2026 contenía tres afirmaciones incorrectas:

\begin{enumerate}
\item Declaraba que la fase comparativa contra el panel experto no se había
      ejecutado por indisponibilidad de evaluadores, y que en consecuencia no
      podían calcularse precisión ni exhaustividad.
\item Reportaba que el detector había marcado 1 de 27 requisitos (3,7\,\%),
      atribuyendo la activación a RF-27 por el término
      \textit{correspondientes}.
\item Incluía entrecomillado un fragmento de RF-27 que no corresponde al texto
      real de ese requisito en el ERS/SRS v2.0.
\end{enumerate}

\textbf{Qué era cierto en realidad, y desde cuándo.} Las tres afirmaciones eran
falsas ya en el momento de escribirse. El panel de tres personas expertas se
ejecutó, sus etiquetas individuales estaban versionadas en
\texttt{06\_Experimento/datos\_procesados/etiquetas\_expertos.csv}, y los
guiones de análisis habían producido \texttt{kappa\_resultados.json},
\texttt{matriz\_confusion\_prf1.json} y \texttt{bootstrap\_ic95.json} con las
cifras correctas. El detector marcó \textbf{0 de 27} requisitos, no 1. El texto
real de RF-27 no contiene el término \textit{correspondientes}.

\textbf{Cómo se detectó.} Mediante auditoría del repositorio contra los
artefactos del pipeline: al ejecutar \texttt{run\_all.py} sobre un clon limpio,
la salida (0/27) no coincidía con la cifra del manuscrito (1/27). La matriz de
confusión publicada en el propio folleto de defensa lo confirmaba de forma
aritmética: con VP\,=\,0 y FP\,=\,0 el detector no puede haber marcado ningún
requisito.

\textbf{Causa raíz.} Las figuras y tablas del manuscrito se habían producido
manualmente, sin ningún guion que las generase a partir de los datos. No
existía, por tanto, ningún mecanismo que hiciera fallar la construcción cuando
el documento y los datos divergían.

\textbf{Mitigación aplicada.}
\begin{enumerate}
\item El manuscrito se reescribió reportando las cifras reales: detector 0/27;
      consenso experto 4/27 (RF-08, RF-17, RF-21, RF-22); matriz de confusión
      VP\,=\,0, FP\,=\,0, FN\,=\,4, VN\,=\,23; precisión, exhaustividad y
      $F_1$ iguales a 0,0000; $\kappa$ de Fleiss 0,2636; $\kappa$ de Cohen por
      pares 0,3478, 0,3571 y 0,0870.
\item Se añadió el guion \texttt{05\_generar\_figuras.py}, que genera todas las
      figuras y todas las tablas de resultados del manuscrito directamente
      desde los artefactos del pipeline. El manuscrito las incorpora con
      \verb|\input|, de modo que ninguna cifra puede ya escribirse a mano sin
      que la discrepancia sea evidente.
\item Se eliminó la cita entrecomillada inexistente.
\end{enumerate}

\textbf{Efecto sobre las conclusiones.} Sustancial y en la dirección correcta.
El resultado real es más informativo que el reportado erróneamente: el detector
no coincide con el consenso experto en ningún caso, y el propio panel muestra
un acuerdo bajo, lo que obliga a discutir simultáneamente el límite del
detector y la estabilidad del criterio de referencia.

\section*{5 bis. DEV-03 --- El corpus \texttt{rf27.json} no coincidía con el
ERS entregado}

\textbf{Qué se registró.} El README del paquete de datos afirmaba que
\texttt{rf27.json} es una salida regenerada de forma determinista desde
\texttt{01\_ERS/ERS\_SRS\_2B\_v2.0.tex}, y que el pipeline es reproducible de
extremo a extremo desde la fuente del ERS.

\textbf{Qué se ejecutó en realidad.} \texttt{rf27.json} se mantenía a mano
desde la Entrega 3 (2A). No recogía las precisiones (umbrales, rangos,
condiciones y aprobaciones de CCB) incorporadas al ERS en la versión 2.0.

\textbf{Momento de detección.} 3 de septiembre de 2026, en auditoría interna
del equipo, que comparó argumento por argumento cada \verb|\rfitem| del .tex
contra el JSON. Resultado: 21 de los 27 requisitos no coincidían.

\textbf{Causa raíz.} Se escribió un extractor determinista
(\texttt{07\_Datos/scripts/extraer\_rf\_desde\_tex.py}) inmediatamente después
de la auditoría, pero no se llegó a ejecutar sobre el corpus real antes de la
entrega evaluada el 13 de septiembre de 2026: el README quedó afirmando una
regeneración que no había ocurrido.

\textbf{Mitigación aplicada, 18/09/2026.} Se ejecutó
\texttt{extraer\_rf\_desde\_tex.py --salida rf27.json} sobre el .tex vigente y
se reejecutó \texttt{run\_all.py} de punta a punta sobre un clon limpio. Las
cifras publicadas (kappa de Fleiss 0,2636; matriz de confusión
VP\,=\,0, FP\,=\,0, FN\,=\,4, VN\,=\,23; bootstrap IC95; análisis de potencia)
son \textbf{idénticas byte a byte} a las obtenidas con el corpus anterior: el
detector congelado se ejecutó sobre ambas versiones y no marcó ningún
requisito en ninguna de las dos.

\textbf{Efecto sobre las conclusiones.} Ninguno. El hallazgo es de veracidad
documental (el README afirmaba una regeneración que no se había hecho), no de
validez de los resultados publicados.

\section*{6. Elementos del protocolo que NO se desviaron}

Se deja constancia expresa de que los siguientes elementos se ejecutaron tal
como fueron registrados, sin cambio alguno:

\begin{itemize}
\item Las tres categorías de patrón del detector y su lista completa de
      expresiones regulares, fijadas antes de conocer las etiquetas expertas.
\item La regla de decisión del detector: un requisito se marca si activa al
      menos una categoría.
\item La regla de consenso del panel: mayoría simple, al menos 2 de 3.
\item El conjunto de métricas: $\kappa$ de Cohen por pares, $\kappa$ de Fleiss,
      matriz de confusión, precisión, exhaustividad y $F_1$.
\item El procedimiento de intervalos de confianza: \textit{bootstrap} no
      paramétrico de 10.000 réplicas con semilla fija igual a 42.
\item El carácter ciego e independiente de la clasificación experta.
\end{itemize}

\section*{7. Trazabilidad}

Este documento se genera desde \texttt{06\_Experimento/osf\_deviations.tex} y
se versiona junto al resto del repositorio. Su contenido se replica en el
apartado \textit{Deviations from pre-registration} del registro
\url{https://osf.io/yp7t3}.

\vspace{10pt}
\noindent
Aprobado por el equipo CarniCore: Castro, Crespo, Gamarra, Pérez y Quintero.

\end{document}
